% TRANSLATION UNIT: translate this file only.
% ===========================================================================
%  03 — Objetivos
%  The \section title is the canonical (Spanish) chapter name; the body prose
%  is written in ENGLISH and translated to Spanish later.
%  Labels OE1..OE5 are kept ready so other chapters can \ref them.
% ===========================================================================

% DONE: One general objective, exactly five specific objectives in the
% QUÉ–CÓMO–PARA–QUÉ shape, a traceability table from each objective to the
% chapter and the instrument that discharges it, the reduction from the ten
% objectives of the proposal to these five, and the claims the objectives
% deliberately do not make.

% EVIDENCE: eight backend services plus one Angular SPA → specs/001-fintcart-platform/plan.md:11
% EVIDENCE: FR-001…FR-031 / SC-001…SC-012 (001), FR-032…FR-085 / SC-013…SC-026 (002),
%   FR-086…FR-127 / SC-027…SC-039 (003) → docs/informe/chapters/06-requerimientos.tex:73-77
% EVIDENCE: twelve constitutional principles, checked as a gate → docs/informe/chapters/05-metodologia.tex:20-27
% EVIDENCE: C4 levels 1 to 3 and the ADR index D-13…D-32 → docs/informe/chapters/07-arquitectura.tex:38-41
% EVIDENCE: 19 primitives, four measured AA ratios, five UI kits, 0 artisanal screens →
%   docs/informe/chapters/09-diseno-ui-ux.tex:10-24; frontend/src/app/shared/ui/index.ts:1-19
% EVIDENCE: ten seed definitions and the formula engine (lexer → parser → AST in JSONB →
%   bounded evaluator in rust_decimal) → services/simulator/src/domain/seeds/mod.rs:144-155;
%   services/simulator/src/domain/formula/mod.rs:1-25
% EVIDENCE: closed block vocabulary validated server-side →
%   services/learning/src/articles/body-doc.validator.ts:1-20
% EVIDENCE: 311/311 unit specs, 59/59 e2e specs, k6 at 200 VUs with 0.00 % HTTP failures,
%   Lighthouse at five widths → docs/informe/chapters/11-pruebas.tex:9-16
% EVIDENCE: four deployment verbs, production image, Caddy edge, 31 migrations across seven
%   databases → docs/informe/chapters/12-despliegue.tex:9-31
% EVIDENCE: 5–8 moderated sessions, ≥ 15 survey responses (target 20), ten-item SUS and a
%   16-item perception survey → docs/informe/04-instrumentos-validacion.md:17-19,282,292-309
% EVIDENCE: preferring is not learning; the instruments cannot demonstrate a learning gain →
%   docs/informe/04-instrumentos-validacion.md:608-616

\section{Objetivos}

\subsection{Objetivo general}

The general objective is one sentence because a project that needs a paragraph to
say what it is for has not yet decided.

\begin{quote}
Design, implement, and evaluate a financial-education platform for the Colombian
context (what), as a polyglot microservice system — eight backend services in Go,
NestJS, and Rust behind a single Angular application, communicating over gRPC
internally and REST only at the edge — built through an incremental,
specification-first process whose gates are executable (how), so that a student in
Colombia can read verified content, practice with graded quizzes, and calculate
with decimal-exact arithmetic on a system whose requirements, behavior, and
usability are each backed by evidence rather than assertion (why).
\end{quote}

The five objectives below decompose that sentence. Each one names three things: the
artefact it produces, the instrument or method by which it produces it, and the end
that makes the artefact worth producing. A reader who cannot point at a chapter and
a measurement for an objective has not been given an objective but a topic.

\subsection{Objetivos específicos}

\begin{itemize}
  \item[OE1.] \label{oe1} Establish the requirements and the architecture of the
  platform (what), through an IEEE Std 830-1998-style specification of 127
  functional requirements and 39 acceptance criteria numbered continuously across
  three feature generations, a C4 description at levels 1 to 3 whose planning
  decisions are promoted to architecture decision records, and a
  database-per-service data model with its migration chain (how), to fix the scope,
  the twelve constitutional principles, and the criterion by which each requirement
  is later accepted or declared not claimed, before implementation begins (why).

  \item[OE2.] \label{oe2} Design the interface and the visual system (what), through
  a token-governed design system of 19 primitives, five full-screen reference kits,
  and a single responsive scale of four breakpoints down to a 360 px floor, audited
  for WCAG 2.1 AA contrast and for a rule that a monetary figure is never truncated
  (how), so that hierarchy is carried by typography and space rather than by
  decoration and a student can locate the title, the amount, and the next action in
  one pass (why).

  \item[OE3.] \label{oe3} Implement the platform's capabilities (what), as eight
  backend services in Go, NestJS, and Rust plus the Angular application, over gRPC
  contracts, with a formula engine that parses, persists the abstract syntax tree
  as \code{JSONB}, and evaluates in \code{rust_decimal}; a content pipeline whose
  block document is validated server-side against a closed vocabulary; and ten seed
  definitions shipped as data (how), so that a student can register, read, practice
  a randomized quiz, run a calculator, and review progress, and so that no monetary
  value can silently lose a digit (why).

  \item[OE4.] \label{oe4} Demonstrate that the implemented system satisfies the
  specification (what), through unit, integration, and end-to-end suites (311 unit
  tests and 59 Playwright specs), a load profile of 200 virtual users, Lighthouse
  runs at five widths, and continuous-integration gates that reject a float
  literal, a breaking contract, or a violation of the report's own glossary (how),
  so that each of the 39 acceptance criteria resolves into reproducible evidence or
  into an explicit statement that it is not claimed (why).

  \item[OE5.] \label{oe5} Put the platform into operation and evaluate it with
  representative users (what), by raising the container topology with four verbs,
  serving a multi-stage image behind a TLS-terminating reverse proxy, and applying a
  protocol of five observed tasks, the ten-item System Usability Scale, and a
  16-item perception survey to 5 to 8 moderated sessions and at least 15 responses
  (how), to establish whether a student completes real tasks unaided and whether the
  interface is trusted with money, reporting that as usability and never as a gain
  in financial literacy (why).
\end{itemize}

\subsection{De diez objetivos a cinco}

The proposal stated ten specific objectives; this report states five, and the
reduction is a decision about what an objective is, not a cut in what was built.
Three of the ten were pairings of one objective with itself. The first two both
establish requirements — the second repeats \enquote{funcionales, no funcionales}
before adding architecture in the same breath — so both are absorbed into OE1,
which carries requirements \emph{and} architecture \emph{and} the data model in one
statement with one closing measurement. The seventh and the eighth both describe
testing, one by naming the levels (unit, integration, end-to-end) and the other by
naming the same integration in different words; integration testing is a step
inside verification, not an objective beside it, and both are absorbed into OE4.
The ninth and the tenth both describe the delivery pipeline, first as
\enquote{ciclos de integración continua} to guarantee delivery and then as
deployment to production through those same cycles.

Four of the ten named an activity instead of a result. \enquote{Utilizando
herramientas de diagrama UML} says which tool the author will open, not what will
exist when the objective closes; \enquote{asegurar el funcionamiento adecuado} and
\enquote{garantizar la entrega} are wishes, because neither names an instrument
that could contradict them. An objective that cannot fail is not an objective, and
an activity is not a deliverable. Six of the ten were steps inside another
objective rather than objectives of their own: the user-experience clause of the
third belongs to design; the event-management components of the fifth are one
capability of the build; the frontend of the sixth is the half of the build that
OE2 designs and OE3 implements.

The absorption, exhaustively: objectives 1 and 2 into OE1; 3 into OE1 for its
architecture and OE2 for its interface clause; 4 and 5 into OE3; 6 into OE2 for the
design and OE3 for the implementation; 7 and 8 into OE4; 9 into OE4, where the
pipeline is the gate rather than the goal; and 10 into OE5, because a platform
nobody can reach cannot be evaluated. Two things were
added that the ten did not have. The first is an evaluation objective with
instruments rather than an intention: the ten could be satisfied by a running
system that no one had ever used. The second is a measurement attached to every
objective, which is what turns a list of wishes into a work plan.

\subsection{Trazabilidad de los objetivos}

Table~\ref{tab:objetivos-trazabilidad} closes the loop. For each objective it names
the chapter where the artefact is reported and the evidence that a reader can
re-run, inspect, or reject. Two rows cite two chapters: the delivery pipeline
belongs to verification as a gate and to evaluation as the environment in which
people meet the product, and the report keeps those two readings apart.

\begin{table}[htbp]
  \centering
  \caption{Each specific objective, the chapter that discharges it, and the evidence that closes it.}
  \label{tab:objetivos-trazabilidad}
  \begin{tabular}{@{}p{1.1cm}p{3.1cm}p{1.7cm}p{8.9cm}@{}}
    \hline
    Obj. & What it establishes & Chapter & Evidence that closes it \\
    \hline
    OE1 & Requirements, architecture, data model & 6, 7, 8 &
    \code{FR-001}--\code{FR-127} and \code{SC-001}--\code{SC-039} in the three
    \code{specs/}; C4 levels 1--3 with the ADR index \code{D-13}--\code{D-32};
    seven PostgreSQL databases and 31 migrations. \\
    OE2 & Interface and visual system & 9 &
    Nineteen primitives exported by \code{shared/ui/index.ts}; four measured AA
    contrast ratios; \code{design-debt.mjs --check} reporting 0 artisanal-class
    screens; five widths, 360 px to 1280 px. \\
    OE3 & Implementation & 10 &
    Eight services plus the SPA; formula engine with exact \code{pot} and
    approximate \code{potd}; ten seed definitions; closed block vocabulary
    validated server-side. \\
    OE4 & Verification and delivery & 11, 12 &
    311 unit tests and 59 e2e specs passing; k6 at 200 virtual users with
    0.00 \% HTTP failures; Lighthouse at five widths; CI jobs for contracts,
    Go, Rust, TypeScript, and the barriers. \\
    OE5 & Operation and evaluation & 12, 13 &
    The four deployment verbs and the production image behind a TLS edge;
    five observed tasks, the ten-item SUS, and 16 survey items, administered to
    5--8 sessions and $\geq 15$ responses. \\
    \hline
  \end{tabular}
\end{table}

\subsection{Lo que los objetivos no afirman}

No objective in this chapter asserts that the platform improves financial literacy.
That claim would require a controlled design — a comparison group, a pre- and
post-measurement, and an instrument validated for the construct — and nothing here
is that. What OE5 measures is whether a student can complete five real tasks
unaided, what the System Usability Scale makes of the interface, and what a
16-item survey records about its perceived legibility, hierarchy, and trust. Those
are properties of the interaction, not of the learning, and chapter~13 reports them
under that name.\todo{chapter~13 is still a stub: the instruments exist and the
protocol is fixed, but no session has been reported, so the evidence for OE5 is
pending rather than absent}. A high usability score is compatible with a platform
that teaches nothing, and a low one with a platform that teaches well; the report
keeps the distinction instead of letting a green number imply the first.

Two smaller refusals follow from the same discipline. OE4's numbers are properties
of the artefact: 311 passing specs say that the code behaves as specified, not that
anyone learns from it, and chapter~11 states that limit where the numbers are
printed. OE5's operational claim is a single-host deployment that serves an address
and terminates TLS, not a highly available one: there is one replica per service
and no backup, and chapter~12 records that as the state of the system rather than
as an aspiration. The load figure of 200 virtual users is a threshold that was met
on one host shared with seven databases and a broker; it is not a capacity plan,
and no objective promises one.
